% Apuntes sobre el paquete beamer: http://metodos.fam.cie.uva.es/~latex/apuntes/apuntes.html (Universidad de Valladolid)
\documentclass{beamer}

\usepackage[utf8x, utf8]{inputenc}
\usepackage{ucs}
\usepackage[T1]{fontenc}
\usepackage{lmodern}
\usepackage[spanish]{babel}
\usepackage{graphicx}


\newcommand{\software}{\textit{software }}

% Configuración de Beamer
\mode<presentation>{\usetheme{Berkeley} \setbeamercovered{dynamic}}
% Otros temas semejantes: Madrid, PaloAlto, Ilmenau, Berkeley, Rochester
\logo{\includegraphics[scale=.04]{logodiedral.pdf}}
\institute{Universidad Complutense de Madrid}

\title[ACE]{\itshape Airline Common Environment}
\author{Grupo $\mathcal{D}$iedral}
\date{20 de diciembre de 2012}

\begin{document}
\maketitle

	\frame{\frametitle{Resumen} \tableofcontents}
	
	\section{Descripción general}
	
\begin{frame}
	\frametitle{Descripción general}
	
	\begin{block}{Resumen}
		\textit{Airline Common Environment} es un producto \software destinado a la gestión integral de una compañía aérea.		\end{block}

	\begin{itemize}
		\item Propósito y alcance
			\begin{itemize}
				\item Gestión interna
				\item Gestión externa
			\end{itemize}
		\item Perspectivas y funciones del producto
		\item Características del usuario
		\item Restricciones y dependencias
	\end{itemize}
\end{frame}

\end{document}
